% ACL 2026-style submission. Compile on Overleaf with pdfLaTeX + BibTeX.
\documentclass[11pt]{article}
\usepackage[review]{acl}
\usepackage{times}
\usepackage{latexsym}
\usepackage[T1]{fontenc}
\usepackage[utf8]{inputenc}
\usepackage{microtype}
\usepackage{inconsolata}
\usepackage{graphicx}
\usepackage{amsmath}
\usepackage{booktabs}
\usepackage{multirow}
\usepackage{enumitem}
\usepackage{xspace}
\usepackage{xcolor}
\usepackage{CJKutf8}

\newcommand{\umr}{\textsc{UMR}\xspace}
\newcommand{\amr}{\textsc{AMR}\xspace}
\newcommand{\sys}{\textsc{umr-critique}\xspace}
\newcommand{\tree}{\textsc{umr-tree}\xspace}
\newcommand{\refine}{\textsc{umr-refine}\xspace}
\newcommand{\skill}[1]{\texttt{#1}}
\newcommand{\zh}[1]{\begin{CJK}{UTF8}{gbsn}#1\end{CJK}}
\newcommand{\todo}[1]{\textcolor{red}{[TODO: #1]}}

\title{Guidelines for Machines, Not for Humans:\\
Self-Evolving Instructions and Tools for Complex Semantic Annotation}

\author{Anonymous ACL submission}

\begin{document}
\maketitle

\begin{abstract}
Annotation guidelines are written for humans: they teach a procedure step by step and illustrate it with many examples. A natural way to use LLMs for annotation is to imitate that procedure. We argue that a guideline for a model should instead describe the target representation, state where the ontology departs from the model's intuitions, and leave the rest to pretraining. Building on instruction learning from gold data \citep{acl2026base}, we extend this view to complex annotation in two ways: models query external ontology knowledge through tools instead of memorizing it from instructions, and a reflection loop can propose new tools and new pipeline steps in addition to rewriting instructions, under machine checks and a review agent. On Chinese Uniform Meaning Representation (\umr) parsing, covering sentence graphs and document-level temporal, modal and coreference relations, our three-call-per-sentence pipeline outperforms a pipeline that imitates the human annotation procedure on four of five held-out documents with the same training data (e.g., 58.5 vs.\ 55.7 sentence F), at about one ninth of the cost. Step-wise ablations show that the transferable gains come from the instructions themselves: dedicated checking steps add 6--9 F on training documents but about 0 F on unseen ones. \todo{one sentence on tool evolution results}
\end{abstract}

%=====================================================================
\section{Introduction}

To carry out or analyze a specific task, practitioners often design a dedicated ontology and annotate data following a guideline. Traditionally this requires large amounts of human labor, and the start-up phase of data production is slow and expensive. Complex semantic annotation is an extreme case. In Uniform Meaning Representation \citep[\umr;][]{bonn2023umr,gysel2021design}, each sentence receives a full semantic graph with concepts, senses, arguments, aspect and named entities, and each document receives an additional layer of temporal, modal and coreference relations. Annotators need long training, and a document can take hours.

Annotation involves two things: identifying the structure in the data, and mapping that structure onto the labels and conventions of the ontology. Through pretraining, large language models (LLMs) have seen enough language to grasp who did what to whom, when, and with what certainty; this half can be left to the model. The other half---PropBank senses and frames \citep{palmer2005propbank}, \umr attribute values, conventions for Chinese constructions---is external knowledge that a model cannot be expected to know.

\citet{acl2026base} \todo{one or two sentences on what the ACL 2026 paper did, e.g., writing guidelines as instructions for a model and rewriting them automatically against gold}. The idea behind this line of work is that \emph{a guideline for a model need not imitate a guideline for humans}. Humans need a procedure and many examples because they learn the representation from scratch and can attend to only a few things at a time. A model already has linguistic and semantic knowledge and can process a whole sentence at once. What it lacks is what the target representation looks like, and where the ontology departs from its intuitions.

Applying this idea to complex annotation such as \umr requires more than rewriting instructions, for two reasons. First, the external knowledge is too large to fit in an instruction: PropBank has thousands of frames, and no instruction can make a model choose sense and argument numbers from memory. Second, some errors are mechanical---an attribute value outside its closed list, an \amr role name used in \umr, an event without aspect---and a short program checks them more reliably than a sentence of advice. In other words, the model needs \textbf{tools} as well as better instructions; and which tools it needs is not known in advance, so they should be learned from errors, just as instructions are.

We extend \citet{acl2026base} in three ways:
\begin{enumerate}[leftmargin=*,itemsep=1pt,topsep=2pt]
  \item \textbf{Tool use.} Every skill can call shared tools, including \texttt{propbank\_lookup} for PropBank frames and \texttt{get\_example} for worked examples. External knowledge is looked up on demand rather than packed into the instruction.
  \item \textbf{Self-evolving tools and skills.} Reflection can not only rewrite instructions but also propose a new tool (a pure Python function callable by every skill) or a new step in the pipeline. To keep evolution under control, every proposal must pass machine checks---an import whitelist, an AST-level ban on file and network access, its own tests run in an isolated process, and no overlap of eight or more characters with any document---and then a strict review agent that asks whether the problem recurs, whether an edit to an existing skill would fix it, whether the new component has one narrow responsibility, whether it is worth an extra call, whether it is correct, and whether it leaks document content.
  \item \textbf{Complex annotation with a short chain.} Each sentence takes three model calls. \skill{generate} writes the whole sentence graph from a description alone, with no examples; \skill{critique} checks it category by category (concepts, senses, attributes, arguments, roles, names, reentrancy) against a list of common errors; a rule step repairs well-formedness; and \skill{document} writes the sentence's temporal, modal and coreference relations in one call. The lists of common errors are accumulated by reflection from gold.
\end{enumerate}

On Chinese \umr, this pipeline outperforms a decision-driven skill orchestration that follows the human annotation procedure step by step \citep[\tree;][]{anon2026orchestration}: with the same training documents, it scores 58.5 vs.\ 55.7 sentence F on 0005, at \$0.08 instead of about \$0.7 per sentence. Step-wise ablations show that the gains come mainly from the instructions: \skill{critique} adds only 0.8 F on average on unseen documents, while reflection raises \skill{generate} alone from 43.5 to 75.7 F on its training document. A variant with five dedicated checking steps (\refine) gains 6--9 F from the checks on its training documents but about 0 F on test documents, suggesting that what extra checking steps learn is largely specific to the documents they were trained on.

Our contributions are:
\begin{itemize}[leftmargin=*,itemsep=1pt,topsep=2pt]
  \item Extending instruction learning for annotation \citep{acl2026base} to tool use, so that ontology knowledge is queried on demand instead of written into instructions.
  \item A mechanism for self-evolving tools and skills, in which reflection proposes new tools and pipeline steps that are admitted only after machine checks and a review agent.
  \item A validation on \umr sentence graphs and document-level relations: three calls per sentence, better than imitating the human procedure, at about one ninth of its cost.
  \item Ablations showing that for complex annotation the transferable gains come from instructions, and that extra checking steps hardly transfer to new documents.
\end{itemize}

%=====================================================================
\section{Related Work}

\paragraph{LLMs for semantic annotation.}
\umr \citep{bonn2023umr,gysel2021design,bonn2025umr2} remains low-resource. \citet{chun2024pipeline} parse documents into \umr via an \amr parser and conversion rules. Direct prompting of LLMs struggles with the formal conventions of \amr and \umr \citep{ettinger2023expert,sun2024chinese}. A companion system \citep{anon2026orchestration} instead decomposes \umr annotation into the human procedure, with the model making local decisions; we use it as our main comparison.

\paragraph{Prompt and instruction optimization.}
Prompts can be optimized by search \citep{zhou2023ape}, by textual ``gradients'' \citep{pryzant2023protegi,yuksekgonul2025textgrad}, by LLMs as optimizers \citep{yang2024opro}, or within compiled pipelines \citep{khattab2024dspy}. \citet{acl2026base} learn annotation instructions from gold \todo{refine}. Like error-driven rule learning \citep{brill1995transformation}, our reflection works from the errors on gold documents, but attributes each error to a responsible skill and an evaluation category, and can change the toolset as well as the text.

\paragraph{Self-correction.}
Self-Refine \citep{madaan2023selfrefine} and Reflexion \citep{shinn2023reflexion} let models revise their own outputs, but intrinsic self-correction without external feedback often fails to help \citep{huang2024cannot}. Our ablations find the same for \umr: a critique step helps little on unseen documents, while feedback from gold, written into the instructions, does transfer.

\paragraph{Tool making.}
Models can learn to call tools \citep{schick2023toolformer}, create tools for a task \citep{cai2024latm,qian2023creator}, and grow a library of skills \citep{wang2023voyager}. We add tool creation to instruction learning for annotation, with an admission process designed to keep the toolset small and safe.

%=====================================================================
\section{Method}

\subsection{Pipeline}

Each sentence is processed by a chain of skills. A skill is an instruction file, optional worked examples, access to tools, and an output schema; the input is placed at the end so that the instruction prefix is cached across calls.

\paragraph{\skill{generate}} receives the sentence and writes its full \umr graph. Its instruction describes what a sentence graph contains (concepts, roles, attributes, named entities, reentrancy) and the format, with no examples. It looks up PropBank frames with \texttt{propbank\_lookup}, and writes \texttt{:modal-strength} on each event.

\paragraph{\skill{critique}} receives the graph and checks it one evaluation category at a time: (a) nodes, (b) senses, (c) attributes, (d) arguments, (e) role labels, (f) named entities, (g) reentrancy. Each category states what the evaluation checks, followed by a list of \emph{common errors}; the seed has only a few generic ones and reflection adds the rest. It returns a list of findings \(\{\textit{check},\textit{problem},\textit{fix}\}\) and the corrected graph.

\paragraph{Rule step.} Without calling a model, the last parseable graph of the chain is repaired: duplicate variables, references to undefined variables, self-loops and parts disconnected from the root are fixed, and unbalanced brackets are corrected.

\paragraph{\skill{document}} receives the sentence's final graph and the five preceding ones and writes all of its temporal, modal and coreference relations at once. Author modality is derived from the \texttt{:modal-strength} on each event.

\subsection{Reflection}

Reflection runs on one training document at a time, for three rounds.

\paragraph{Find.} For each sentence, a reflection model compares gold and prediction (with AnCast's variable alignment and the document-level relations) and lists every problem, labeled with its evaluation category, its origin (the \skill{generate} draft or a \skill{critique} change), and the skill that should be changed: \skill{critique} if it could have caught the error in that category, \skill{generate} if the draft misses something \skill{critique} cannot repair, \skill{document} for document-level relations.

\paragraph{Plan and apply.} Problems are grouped by skill (and by category within \skill{critique}), and the reflection model proposes \emph{modify}, \emph{add} or \emph{delete} edits to instructions and examples; common errors go into the list of the corresponding category. Edits are applied to a copy of the skill set, which becomes the next version. Any edit that shares eight or more consecutive characters with any document is rejected.

\subsection{Evolving tools and skills}

In each round, reflection may also propose at most one new tool and one new skill, and edits may declare that they require it.

\paragraph{New tools} are pure Python functions with type annotations and a docstring, callable by every skill. They may import only \texttt{re}, \texttt{json}, \texttt{math}, \texttt{itertools}, \texttt{collections}, \texttt{functools}, \texttt{unicodedata}, \texttt{string} and \texttt{penman}. An AST check rejects \texttt{open}, \texttt{eval}, \texttt{exec}, \texttt{getattr} and access to underscored attributes, and requires exactly one public function. The proposal must include at least three tests, which run in an isolated subprocess and must all pass.

\paragraph{New skills} are additional steps in the chain with the same interface as \skill{critique} (graph in, graph and changes out; or relations in, relations and changes out). The registry must load, the link must be valid, and the instruction prefix must be identical across calls.

\paragraph{Review.} Proposals that pass the machine checks go to a review agent with six criteria: (1) the problem recurs and is costly; (2) it cannot be fixed by editing an existing skill; (3) the component has one narrow responsibility and does not conflict with other skills; (4) it is worth an extra call; (5) its content is correct; (6) it leaks no document content. The agent approves only if all hold, and rejects when unsure. Edits that require a rejected component are dropped.

%=====================================================================
\section{Experimental Setup}

\paragraph{Data and protocol.}
We use Chinese \umr documents 0001--0007 after the cleanup of \citet{chun2024pipeline}. Training documents are processed in order, three reflection rounds each; version $R_{xy}$ is the skill set after round $y$ on document $x$. The first parse of each document uses the version trained only on earlier documents, so it is a test on an unseen document, and the sequence of first parses is a learning curve.

\paragraph{Systems.}
\textbf{\sys} is our pipeline. \textbf{\refine} replaces \skill{critique} with five checking skills (concepts, arguments, attributes, entities, reentrancy) and has the same reflection machinery (it may add a tool or skill at most once every three documents); it takes seven calls per sentence. \textbf{\tree} \citep{anon2026orchestration} follows the human procedure with dozens of decisions per sentence and reflection confined to the skill that made each decision. \textbf{Single call}: one call per sentence with a minimal prompt and no document level. All three skill-based systems share the same conventions in their seeds.

\paragraph{Models and metrics.}
Parsing uses Claude Sonnet 5.5 and reflection Claude Opus 5.5. We report AnCast++ \citep{sun2024ancast} sentence F and the comprehensive score, and use a fine-grained judge for per-category analysis.

%=====================================================================
\section{Results}

\begin{table}[t]
\centering\small
\setlength{\tabcolsep}{3pt}
\begin{tabular}{@{}lccc@{}}
\toprule
Doc & \tree & \refine & \sys \\
\midrule
0001 (seed) & 40.8 / 38.4 & 46.0 / 41.8 & 45.1 / 40.5 \\
0002 ($R_{13}$) & 59.0 / 54.4 & 59.7 / 55.3 & \textbf{61.1} / \textbf{56.2} \\
0003 ($R_{23}$) & \textbf{57.9} / \textbf{56.6} & -- & 57.2 / 55.4 \\
0004 & 50.0 / 48.5 & 51.9 / \textbf{50.0} & \textbf{52.4} / 49.7 \\
0005 & 55.7 / 54.3 & 57.5 / 55.5 & \textbf{58.5} / \textbf{57.1} \\
0006 & 49.5 / 45.4 & -- & \textbf{56.1}$^\dagger$ / \textbf{49.2}$^\dagger$ \\
\midrule
calls / sent. & dozens & 7 & 3 \\
\$ / sent. & $\sim$0.71 & $\sim$0.16 & $\sim$0.08 \\
\bottomrule
\end{tabular}
\caption{Sentence F / comprehensive on documents each version has never seen. Versions: 0004 uses \tree $R_{32}$, \refine $R_{23}$, \sys $R_{33}$; 0005 uses \tree $R_{42}$, \refine $R_{23}$, \sys $R_{43}$; 0006 uses \tree $R_{42}$ and \sys $R_{53}$ ($^\dagger$also trained on 0005). Cost at API prices. \todo{run \sys $R_{43}$ on 0006 and 0007 for a matched comparison}}
\label{tab:main}
\end{table}

\paragraph{Main comparison.}
Table~\ref{tab:main} compares the three systems on documents each version has not seen. With the same training documents, \sys matches or beats \tree on every document except 0003 (0.7 F lower), and is 2.4 F higher on 0004 and 2.8 F higher on 0005, with 3 calls per sentence instead of dozens. \refine, trained on only two documents, is close to \sys. All three are far above single-call generation (32.9 F on 0004+0005).

\paragraph{Learning on training documents.}
Within a document, reflection rounds raise sentence F substantially (Table~\ref{tab:train}): from 45.1 to 77.5 on 0001 and from 57.2 to 84.4 on 0003. Each round applies 17--64 edits; most go to \skill{generate} and \skill{critique} (3--34 and 8--27 per round), and 2--7 to \skill{document}. The number of problems found on a document falls across rounds (0001: 752 $\rightarrow$ 516; 0003: 407 $\rightarrow$ 250).

\begin{table}[t]
\centering\small
\setlength{\tabcolsep}{4pt}
\begin{tabular}{@{}lcccc@{}}
\toprule
Doc & iter 1 & iter 2 & iter 3 & after \\
\midrule
0001 & 45.1 & 64.4 & 74.2 & 77.5 \\
0002 & 61.1 & 68.5 & 76.5 & 77.3 \\
0003 & 57.2 & 72.9 & 83.6 & 84.4 \\
0004 & 52.4 & 62.9 & 70.0 & 72.4 \\
0005 & 58.5 & 69.6 & 74.0 & 79.7 \\
0006 & 56.1 & 74.0 & 79.5 & 81.0 \\
\bottomrule
\end{tabular}
\caption{\sys sentence F on training documents. Iteration 1 parses with the version trained on earlier documents only.}
\label{tab:train}
\end{table}

%=====================================================================
\section{Analysis}

\subsection{Where the gains come from}

We score the sentence graph after each step of the chain (without the document layer).

\paragraph{\skill{critique} adds little.}
On unseen documents (Table~\ref{tab:steps}), \skill{critique} changes sentence F by $-0.4$ to $+1.6$, 0.8 on average. Reflection, meanwhile, improves the \skill{generate} draft itself: on 0001, the draft rises from 43.5 (seed) to 62.1, 70.2 and 75.7 over three rounds. Most of what reflection learns ends up making the first draft right.

\begin{table}[t]
\centering\small
\setlength{\tabcolsep}{4pt}
\begin{tabular}{@{}llccc@{}}
\toprule
Doc & Version & \skill{generate} & +\skill{critique} & $\Delta$ \\
\midrule
0001 & $R_0$ & 43.5 & 45.1 & +1.6 \\
0002 & $R_{13}$ & 61.5 & 61.1 & $-$0.4 \\
0003 & $R_{23}$ & 56.6 & 57.2 & +0.6 \\
0004 & $R_{33}$ & 51.0 & 52.4 & +1.4 \\
0005 & $R_{43}$ & 57.9 & 58.5 & +0.6 \\
0006 & $R_{53}$ & 55.4 & 56.1 & +0.7 \\
\bottomrule
\end{tabular}
\caption{Sentence F after each step of \sys on unseen documents.}
\label{tab:steps}
\end{table}

\paragraph{Extra checking steps do not transfer.}
\refine has five checking skills after \skill{generate}. Table~\ref{tab:refine} contrasts their combined contribution on training documents (parsed after the document's own rounds) and on test documents: +9.2 and +8.9 F on the training documents, but between $-1.3$ and $+1.5$ on test documents, 0.0 on average. Reflection writes fixes into the checks that work on the document it saw; what transfers to new documents is what goes into \skill{generate}'s instruction.

\begin{table}[t]
\centering\small
\setlength{\tabcolsep}{4pt}
\begin{tabular}{@{}lllccc@{}}
\toprule
& Doc & Ver. & gen. & +5 checks & $\Delta$ \\
\midrule
\multirow{2}{*}{train} & 0001 & $R_{13}$ & 65.8 & 75.0 & +9.2 \\
 & 0002 & $R_{23}$ & 67.0 & 75.9 & +8.9 \\
\midrule
\multirow{4}{*}{test} & 0004 & $R_{13}$ & 49.1 & 50.0 & +0.9 \\
 & 0004 & $R_{23}$ & 50.4 & 51.9 & +1.5 \\
 & 0005 & $R_{13}$ & 58.9 & 57.6 & $-$1.3 \\
 & 0005 & $R_{23}$ & 58.6 & 57.5 & $-$1.1 \\
\bottomrule
\end{tabular}
\caption{\refine: sentence F after \skill{generate} and after its five checking skills.}
\label{tab:refine}
\end{table}

\subsection{Tool and skill evolution}

Over 20 rounds (three on each of 0001--0006, two on 0007), reflection proposed three tools and no skills, and the review agent rejected all three. \texttt{check\_umr\_graph} (round 1 on 0001) checked colon-prefixed attribute values, the closed list of \texttt{:modal-strength} values, \texttt{person} vs.\ \texttt{individual-person} and \amr roles; the reviewer judged the colon and closed-list checks worth keeping but found that about half of the checks could be fixed by editing instructions. \texttt{check\_umr\_form} (round 2) checked bare constants and aspect on connectives; the reviewer accepted the first idea but found that the second contradicted \skill{generate}'s own instruction and example. \texttt{check\_doc\_relations} (round 3 on 0003) checked coreference on name nodes and events with two temporal parents; the reviewer agreed the errors recur and suit a deterministic check, but found part of its logic incorrect. In each case the reviewer separated a sound core from a flawed rest. \todo{either (a) let rejected proposals be revised in the next round with the review attached, and report an accepted tool with an ablation; or (b) present controlled rejection as the result}

%=====================================================================
\section{Conclusion}

A guideline for a model should not imitate a guideline for humans. For complex semantic annotation, a short chain whose first step writes the whole graph from a description, with tools for external knowledge and instructions learned from gold, outperforms a pipeline that follows the human procedure, at a fraction of the cost. The transferable gains come from the instructions; extra checking steps mostly learn document-specific fixes. Letting reflection evolve tools as well as text is promising, and the admission process keeps it controlled.

\section*{Limitations}

Experiments are on Chinese \umr only, with seven evaluated documents. The comparison in Table~\ref{tab:main} uses versions trained on the same documents but with different numbers of rounds (\tree used two rounds on 0003 and 0004, \sys three), and \sys's version for 0006 had also seen 0005. No tool or skill has yet been admitted, so the effect of evolved tools is not measured. Each configuration was run once. Costs are at API prices for one model family.

\bibliography{references}

\end{document}
